% -*- TeX-master: "../main.tex" -*-
\chapter{Makieta}
Kod źródłowy makiety \parencite{makieta} pod adresem \parencite{repo}. Została ona stworzona w czystym html'u, css'ie oraz javascript'cie bez używania żadnych frameworków ani AI w celu zachowania pełnej kontroli nad wyglądem makiety oraz animacjami.



\section{Symulacja przeciągania palcem}
Makieta pozwala na wykonywanie gestów przeciągnięcia myszką których obsługa została zaimplementowana w minimalnej funkcji.


\begin{lstlisting}[language=javascript, caption={Kod odpowiadający za gest przeciągania.}, label={code:przeciagniecie}]
function attachSwipeHandler(el=document.getElementById(""),fn=()=>{},right=true,corner=5,sens=3,updown=false){
    var lastP = 0;
    var eligible = false;
    el.addEventListener("mousedown" ,(/** @type {MouseEvent} **/ev)=>{
        if(lock)
            return;
        
        let p;
        if(!updown)
            p = ev.layerX*100/el.getBoundingClientRect().width;
        else
            p = ev.layerY*100/el.getBoundingClientRect().height;

        if(!right) p = 100-p;

        // console.log("md "+el.tagName+" "+Math.round(p));

        if(p<corner){
            // debugger
            console.log("eligible "+p+" < "+corner)
            eligible=true;
            lastP = p;
        }else eligible=false
        
    });
    el.addEventListener("mousemove" ,(/** @type {MouseEvent} **/ev)=>{
        if(lock)
            return;
        
        if(!updown)
            p = ev.layerX*100/el.getBoundingClientRect().width;
        else
            p = ev.layerY*100/el.getBoundingClientRect().height;

        if(!right) p = 100-p;
        
        // console.log("mm "+el.tagName+" "+Math.round(p));

        if(eligible){
            if(lastP+sens<=p){
                fn();
                eligible = false;
                console.log("triggered");
            }
        }
    });
    el.addEventListener("mouseup" ,(/** @type {MouseEvent} **/ev)=>{
        lock = false;
        // console.log("mu "+el.tagName+" ");
        eligible = false;
    });
}
\end{lstlisting}

Jak widać w kodzie \ref{code:przeciagniecie} funkcja działa wyłącznie dla myszki ponieważ w dla przeglądarki internetowej dotyk ekranu jest osobnym zbiorem eventów (to jest "touchstart" zamiast "mousedown") który wymaga wsparcia dla tak zwanego multitouch (po polsku wielo dotyku) co utrudnia prostą implementacje. Z tego powodu na laptopach z ekranami dotykowymi oraz na telefonach symulacja przeciągnięcia nie działa.
\newline\newline
Funkcje można wykorzystać do chowania slider'ów oraz jumper'ów zamiast klikania w przyciemnioną część ekranu. Dodatkowo pozwala również ona na przeciągnięcie od boku ekranu w edytorze kodu żeby wysunąć slider (rysunek \ref{fig:code_screen}) z listą plików w projekcie (rysunek \ref{fig:code_doc_screen}) bez potrzeby klikania przycisku hamburgera w górnym rogu ekranu.



\section{Symulacja wibracji urządzenia}
Makieta pozwala na symulacje wibracji telefonu po przez potrząśnięcie ekranu w momencie próby zatwierdzenia akcji bez wpisania informacji w polu tekstowym, funkcja ta jest wykorzystywana przy tworzeniu nowego projektu (rysunek \ref{fig:new_proj_screen}), dodawaniu repozytorium git (rysunek \ref{fig:import_git_screen}) oraz przy tworzeniu commit'a (rysunek \ref{fig:git_screen})

Kod który za to odpowiada jest jeszcze prostszy niż kot odpowiadający za przeciąganie.

\begin{lstlisting}[language=javascript, caption={Kod odpowiadający za wywoływania wibracji dodawanie czerwonej ramki do pustych pól.}, label={code:shake}]
function validateText(/**@type {Element}**/inp,fn=()=>{}){
    /**@type {string}**/
    let txt = inp.value
    
    if(txt==""){
        addTag(inp,"inp_wrong");
        addTag(app_screen,"shake")
        setTimeout(()=>{rmTag(app_screen,"shake")},400)
    }else{
        rmTag(inp,"inp_wrong");
        fn();
    }
}
\end{lstlisting}

\begin{lstlisting}[language=css, caption={Kod css odpowiadający za trzęsienie.}, label={code:shake-css}]
@keyframes shakeAnimation {
  from {transform: translate(-1vh,0);}
  to {  transform: translate(1vh,0);}
}

.shake{
  animation-name: shakeAnimation;
  animation-duration: 100ms;
  animation-iteration-count: infinite;
}
\end{lstlisting}

\section{Symulacja przytrzymania elementu}
\label{sec:hold_sym}
Makieta posiada również możliwość symulacji przytrzymania palcem elementu po przez kliknięcie i przytrzymanie klawisza myszki, jeśli klawisz zostanie podniesiony lub usunięty z guzika przed upływem odpowiedniej ilości czasu akcja zostanie anulowana.

Kod odpowiadający za tą funkcje załączam poniżej.

\begin{lstlisting}[language=javascript, caption={Kod odpowiadający obsługę przytrzymanie klawiszu myszki na elemencie.}, label={code:hold}]
function attachHoldHandler(el=document.getElementById(""),fn,wait=500){//500 milisekund
    var clicked = false;
    var id = 0;
    el.addEventListener("mousedown" ,(/** @type {MouseEvent} **/ev)=>{
        let remid = id;
        clicked = true;
        setTimeout(()=>{
            console.log(remid+" x "+id)
            if(clicked&&id==remid){
                fn()
            }
        },wait)
    })
    el.addEventListener("mouseleave",()=>{
        if(clicked){
            id++;
            clicked = false;
        }
    })
    el.addEventListener("mouseup",()=>{
        if(clicked){
            id++;
            clicked = false;
        }
    })
}
\end{lstlisting}

Jak widać w kodzie \ref{code:hold} przytrzymanie również można wywołać jedynie z poziomu myszki.


